Даны две матрицы размерами $M \times N$. Заполнить новую матрицу по следующему правилу: чётные столбцы из первой матрицы, нечётные --- из второй.
Даны две матрицы размерами $M \times N$. Заполнить новую матрицу по следующему правилу: чётные строки из первой матрицы, нечётные --- из второй.
Даны две матрицы размерами $M \times N$. Заполнить новую матрицу по следующему правилу: нечётные столбцы из первой матрицы, чётные --- из второй.
Даны две матрицы размерами $M \times N$. Заполнить новую матрицу по следующему правилу: нечётные строки из первой матрицы, чётные --- из второй.
Даны две матрицы размерами $M \times N$. Заполнить новую матрицу по следующему правилу: столбцы с номерами, кратными трём, из первой матрицы, остальные --- из второй.
Даны две матрицы размерами $M \times N$. Заполнить новую матрицу по следующему правилу: строки с номерами, кратными трём, из первой матрицы, остальные --- из второй.
Даны две матрицы размерами $M \times N$. Заполнить новую матрицу по следующему правилу: первая треть столбцов из первой матрицы, остальные --- из второй.
Даны две матрицы размерами $M \times N$. Заполнить новую матрицу по следующему правилу: первая треть строк из первой матрицы, остальные --- из второй.
Даны две матрицы размерами $M \times N$. Заполнить новую матрицу по следующему правилу: центральная треть столбцов из первой матрицы, остальные --- из второй.
Даны две матрицы размерами $M \times N$. Заполнить новую матрицу по следующему правилу: центральная треть строк из первой матрицы, остальные --- из второй.
Даны две матрицы размерами $M \times N$. Заполнить новую матрицу по следующему правилу: последняя треть столбцов из первой матрицы, остальные --- из второй.
Даны две матрицы размерами $M \times N$. Заполнить новую матрицу по следующему правилу: последняя треть строк из первой матрицы, остальные --- из второй.
Даны две матрицы размерами $M \times N$. Заполнить новую матрицу по следующему правилу: поэлементная сумма двух матриц.
Даны две матрицы размерами $M \times N$. Заполнить новую матрицу по следующему правилу: поэлементная разность первой и второй матриц.
Даны две матрицы размерами $M \times N$. Заполнить новую матрицу по следующему правилу: поэлементное произведение двух матриц.
Даны две матрицы размерами $M \times N$. Заполнить новую матрицу по следующему правилу: поэлементное среднее арифметическое двух матриц.
Даны две матрицы размерами $M \times N$. Заполнить новую матрицу по следующему правилу: модуль поэлементной разности первой и второй матриц.
Даны две матрицы размерами $M \times N$. Заполнить новую матрицу по следующему правилу: поэлементная сумма квадратов элементов двух матриц.
Даны две матрицы размерами $M \times N$. Заполнить новую матрицу по следующему правилу: результат деления элементов первой матрицы на соответствующие элементы второй (если элемент второй не ноль, иначе ноль).
Даны две матрицы размерами $M \times N$. Заполнить новую матрицу по следующему правилу: поэлементное максимум из двух матриц.
Даны две матрицы размерами $M \times N$. Заполнить новую матрицу по следующему правилу: поэлементный минимум из двух матриц.
Даны две матрицы размерами $M \times N$. Заполнить новую матрицу по следующему правилу: поэлементное сравнение, единица, если элемент первой матрицы больше соответствующего элемента второй, иначе ноль.
Даны две матрицы размерами $M \times N$. Заполнить новую матрицу по следующему правилу: поэлементное сравнение, двойка, если оба элемента чётные; единица, если один из элементов чётный, иначе ноль.
Даны две матрицы размерами $M \times N$. Заполнить новую матрицу по следующему правилу: все элементы из первой матрицы, кроме последнего столбца ---его взять из второй.
Даны две матрицы размерами $M \times N$. Заполнить новую матрицу по следующему правилу: все элементы из первой матрицы, кроме последней строки ---её взять из второй.
Даны две матрицы размерами $M \times N$. Заполнить новую матрицу по следующему правилу: все элементы из второй матрицы, кроме первого столбца ---его взять из первой.
Даны две матрицы размерами $M \times N$. Заполнить новую матрицу по следующему правилу: все элементы из второй матрицы, кроме первой строки ---её взять из первой.
Даны две матрицы размерами $M \times N$. Заполнить новую матрицу по следующему правилу: поэлементное произведение первой матрицы на число 2 плюс вторая матрица.
Даны две матрицы размерами $M \times N$. Заполнить новую матрицу по следующему правилу: из первой матрицы взять элементы, если они больше нуля, иначе ---из второй.
Даны две матрицы размерами $M \times N$. Заполнить новую матрицу по следующему правилу: из первой матрицы взять элементы, если они нечётные, иначе ---из второй.
